\documentclass[10pt,a4paper]{amsart}
\usepackage[margin=19mm]{geometry}
\usepackage{amsmath,amssymb,mathtools}
\usepackage[scr=rsfs,cal=euler]{mathalfa}
\usepackage{xfrac}
\usepackage[unicode,hidelinks,bookmarks=false]{hyperref}
\newcommand{\ie}{i.e.\ }
\newcommand{\cf}{cf.\ }
\newcommand{\resp}{respectively\ }
\newcommand{\Subsection}{\subsection}
\providecommand{\nobreakdash}{\nobreak\hbox{-}}
\setlength{\emergencystretch}{2em}
\multlinegap0pt
\usepackage{bm}
\usepackage{bbm}
\usepackage{dsfont}
\usepackage{xspace}
\usepackage{cancel}
\usepackage{mathtools}
\usepackage{amsthm}
\makeatletter
\@ifundefined{theorem}{\newtheorem{theorem}{Theorem}}{}
\@ifundefined{lemma}{\newtheorem{lemma}[theorem]{Lemma}}{}
\makeatother
\newcommand*{\ExtGraph}{\mathcal{E}}
\newcommand*{\LExt}{\operatorname{E}^-}
\newcommand*{\RExt}{\operatorname{E}^+} 
\newcommand*{\cA}{\mathcal{A}}
\newcommand*{\cB}{\mathcal{B}}
\newcommand*{\emptyw}{\varepsilon}
\newcommand*{\from}{\colon}
\newcommand*{\lang}{\mathcal{L}}
\newcommand*{\ret}{\mathcal{R}}
\title[Algebraic characterization of dendricity]{Algebraic characterization of dendricity}
\author{France Gheeraert, Herman Goulet-Ouellet, Julien Leroy, Pierre Stas}
\date{Source: arXiv:2406.15075v1 (CC BY 4.0)}
\begin{document}
\maketitle

\noindent\emph{Source and licence.} Algebraic characterization of dendricity. Version arXiv:2406.15075v1 (CC BY 4.0), distributed under the Creative Commons Attribution 4.0 International licence, as recorded in the arXiv licence metadata for this exact version and shown on the abstract page https://arxiv.org/abs/2406.15075v1. Full rights notice: licenses/arxiv-2406.15075-CC-BY-4.0.txt.\par\smallskip

\noindent\textit{Editorial scope.} Counted unit: the Lemma whose source label is \texttt{L:isomorphism generalized extension graph} in the article (source0.tex lines 159--161, source label \texttt{L:isomorphism generalized extension graph}), delivered together with the article's own complete original proof of it (source0.tex lines 163--178, one \texttt{proof} environment of 16 lines). No mathematics is written, repaired or re-derived: statement and proof are the article's own text line for line. Required context retained uncounted: none. Every definition, display and label of those lines is retained unchanged. Omitted from this package and not redistributed: 1--158, 162--162, 179--264. Nothing inside a retained excerpt is omitted or altered: every retained body is delivered exactly as the article's lines are written, so no omission receipt of any kind accompanies this package. Citations: the retained excerpts carry no citation command at all, so no bibliography entry of the article is transcribed and no reference is redistributed. Reference adaptation: none was needed and none was made; every reference inside the delivered unit resolves inside it, so no unresolved reference or citation is produced. Printed slips kept literal: no printed word, symbol, hypothesis or number of the counted statement or of its proof was altered. Layout and class: the article's own class is \texttt{amsart} with packages including xcolor, amsmath, amssymb, amsthm, doi, csquotes, booktabs, enumitem, graphicx, microtype, hyperref, tikz, cleveref, natbib; the wrapper is the lane's pinned \texttt{amsart} editorial wrapper at 10pt on A4, which restates the theorem-like environments the retained text needs and adds the article's own macro definitions that the retained text needs (\texttt{\textbackslash ExtGraph}, \texttt{\textbackslash LExt}, \texttt{\textbackslash RExt}, \texttt{\textbackslash cA}, \texttt{\textbackslash cB}, \texttt{\textbackslash emptyw}, \texttt{\textbackslash from}, \texttt{\textbackslash lang}, \texttt{\textbackslash ret}), with extra wrapper packages (bm, bbm, dsfont, xspace, cancel, mathtools). The delivered numbering is the unit's own. Assets: no retained span contains a figure environment, a graphics-inclusion command, a PDF-page inclusion, a table or a TikZ picture; no image file is copied into this package, the package declares no asset dependency and no image file is redistributed with it. Licence: version arXiv:2406.15075v1 of this article is distributed under the Creative Commons Attribution 4.0 International licence, as recorded in the arXiv licence metadata for this exact version and shown on the abstract page https://arxiv.org/abs/2406.15075v1. A full rights notice is delivered at licenses/arxiv-2406.15075-CC-BY-4.0.txt. Characters: every retained excerpt is pure ASCII, so the delivered engine, XeLaTeX with its default Latin Modern text font, reports no missing character.
\begin{lemma}\label{L:isomorphism generalized extension graph}
    Let $X$ be a minimal shift space and $w \in \lang(X)$. Then $\ExtGraph_X(w)$ is the image of $\ExtGraph_{D_w(X)}(\varepsilon)$ under a graph morphism.
\end{lemma}

\begin{proof}
    Let $\theta_w\from \cB\to\ret_X(w)$ be the derivating substitution used to define $D_w(X)$. We consider the set 
    \[
        \ret'_X(w) = \{r' \in \cA^+ \mid wr'\in \lang(X)\cap \cA^*w\setminus \cA^+w\cA^+\}.
    \]
    of \emph{right return words} to $w$ and we consider $\theta'_w\from \cB^* \to \cA^*$ defined by $\theta'_w(u) = w^{-1}\theta_w(u)w$ (thus it is a bijection between $\cB$ and $\ret_X'(w)$). Observe that $z\in\lang(D_w(X))$ if and only if $\theta_w(z)w\in\lang(X)$, if and only if $w\theta_w'(z)\in\lang(X)$.
    
    Consider the map $\Theta$ defined on the vertices of $\ExtGraph_{D_w(X)}(\varepsilon)$ as follows:
    \begin{itemize}
        \item $r\in\LExt_{D_w(X)}(\varepsilon)$ is mapped to the last letter of $\theta_w(r)$;
        \item $s\in\RExt_{D_w(X)}(\varepsilon)$ is mapped to the first letter of $\theta_w'(s)$.
    \end{itemize}
    Assume that $r\in \LExt_{D_w(X)}(\emptyw)$ and $s\in \RExt_{D_w(X)}(\emptyw)$ are connected by an edge in $\ExtGraph_{D_w(X)}(\emptyw)$; in other words, $rs\in\lang(D_w(X))$. It then follows that $\theta_w(rs)w = \theta_w(r)w\theta_w'(s)\in\lang(X)$, and thus $\Theta(r)w\Theta(s)\in\lang(X)$. This shows that $\Theta$ defines a graph morphism from $\ExtGraph_{D_w(X)}(\emptyw)$ to a subgraph of $\ExtGraph_X(w)$. 
    
    To show that it is onto, assume that $a\in\LExt_X(w)$ and $b\in\RExt_X(w)$ are connected by an edge in $\ExtGraph_{X}(w)$; in other words $awb\in\lang(X)$. By uniform recurrence, there exist words $u$ and $v$ such that $uw$ starts with $w$ and ends with $aw$, $wv$ starts with $wb$ and ends with $w$, and $uwv\in\lang(X)$. Assuming that $u$ and $v$ are of minimal lengths, we get $u\in\ret_X(w)$ and $v\in\ret_X'(w)$, hence there exist $r,s\in\cB$ such that $\theta_w(r)=u$ and $\theta_w'(s)=v$. Moreover, $\theta_w(r)w\theta_w'(s) = \theta_w(rs)w \in \lang(X)$ so $rs \in \lang(D_w(X))$. Considering $r$ as an element of $\LExt_{D_w(X)}(\emptyw)$ and $s$ as an element of $\RExt_{D_w(X)}(\emptyw)$, $r$ and $s$ are connected in $\ExtGraph_{D_w(X)}(\emptyw)$ and satisfy $\Theta(r)=a$ and $\Theta(s)=b$, which proves that $\Theta$ is onto.
\end{proof}

\end{document}
